% Build with lualatex ai-skills-for-reliable-software.tex (run twice for references).
% Requires Beamer, TikZ, fontspec, Liberation Sans, and DejaVu Sans Mono.
% For a presenter copy, replace "hide notes" with "show notes" and rebuild.
\documentclass[aspectratio=169,10pt]{beamer}
\usepackage{fontspec}
\usepackage{tikz}
\setsansfont{Liberation Sans}
\setmonofont{DejaVu Sans Mono}[Scale=MatchLowercase]
\setbeamertemplate{navigation symbols}{}
\setbeameroption{hide notes}
% \setbeameroption{show notes}
\setbeamerfont{note page}{size=\tiny}
\title{AI skills for reliable software}
\subtitle{Evidence from the mathematical-abstractions repository}

% Slide coordinates use a 13.333 by 7.5 grid, with the origin at the top left.
\newcommand{\SlidePanel}[5]{%
  \fill[fill=#5] (#1,#2) rectangle ++(#3,#4);%
}
\newcommand{\SlideText}[7]{%
  \node[anchor=north west,inner sep=0pt,outer sep=0pt,text=#5,align=left]
    at (#1,#2) {\begin{minipage}[t]{#3\SlideUnit}%
    \raggedright\fontsize{#4pt}{\fpeval{#4*1.15}pt}\selectfont
    #6#7\end{minipage}};%
}
\newlength{\SlideUnit}
\setlength{\SlideUnit}{12mm}
\newenvironment{slidecanvas}{%
  \begin{tikzpicture}[remember picture,overlay,x=\SlideUnit,y=-\SlideUnit]
  \begin{scope}[shift={(current page.north west)}]%
}{\end{scope}\end{tikzpicture}}
% The Beamer section number selects the active step in the roadmap.
\newcommand{\SlideProgress}{%
  \foreach \step/\labeltext in {1/Foundations,2/Scheduler,3/Wet grass,4/Takeaways}{%
    \pgfmathsetmacro{\stepx}{.72+(\step-1)*2.99}%
    \ifnum\insertsectionnumber=\step\relax
      \SlidePanel{\stepx}{.14}{2.91}{.34}{teal}%
      \SlideText{\fpeval{\stepx+.09}}{.205}{2.73}{7.5}{white}{\bfseries}{\step\enspace \labeltext}%
    \else
      \SlidePanel{\stepx}{.14}{2.91}{.34}{panel}%
      \SlideText{\fpeval{\stepx+.09}}{.205}{2.73}{7.5}{muted}{}{\step\enspace \labeltext}%
    \fi
  }%
}
\newcommand{\SlideHeader}[4]{%
  \ifnum\insertsectionnumber>0\relax\SlideProgress\fi
  \SlideHeading{#1}{#2}{#3}{#4}%
}
\newcommand{\SlideHeading}[4]{%
  \SlidePanel{.72}{.66}{.34}{.06}{#3}
  \SlideText{1.2}{.56}{11}{5.5}{#3}{\bfseries}{#2}
  \SlideText{.72}{.91}{11.92}{18.5}{#4}{\bfseries}{#1}
}
\newcommand{\SlideFooter}[3]{%
  \SlidePanel{.72}{7.02}{11.88}{.018}{#3}
  \SlideText{.73}{7.15}{11.35}{5}{#2}{}{#1}
  \SlideText{11.85}{7.12}{.75}{5.5}{#2}{\bfseries}{\raggedleft\insertframenumber/\inserttotalframenumber}
}
\newcommand{\SlideCard}[6]{%
  \SlidePanel{#1}{#2}{#3}{#4}{panel}
  \SlideText{\fpeval{#1+.23}}{\fpeval{#2+.19}}{\fpeval{#3-.46}}{10.5}{teal}{\bfseries}{#5}
  \SlideText{\fpeval{#1+.23}}{\fpeval{#2+.79}}{\fpeval{#3-.46}}{9.5}{ink}{}{#6}
}
\newcommand{\SlideCode}[7]{%
  \SlidePanel{#1}{#2}{#3}{#4}{dark}
  \SlideText{\fpeval{#1+.22}}{\fpeval{#2+.19}}{\fpeval{#3-.44}}{7}{mint}{\bfseries}{#5}
  \SlideText{\fpeval{#1+.22}}{\fpeval{#2+.67}}{\fpeval{#3-.44}}{#6}{light}{\ttfamily}{#7}
}

\definecolor{ink}{HTML}{16302F}
\definecolor{paper}{HTML}{F4F5EE}
\definecolor{teal}{HTML}{007767}
\definecolor{muted}{HTML}{466362}
\definecolor{rule}{HTML}{DDE6DB}
\definecolor{dark}{HTML}{112C2A}
\definecolor{light}{HTML}{F1F7EF}
\definecolor{mint}{HTML}{A9EAD0}
\definecolor{darkrule}{HTML}{48615D}
\definecolor{panel}{HTML}{E7EEE3}
\definecolor{alternate}{HTML}{EDF1E8}
\definecolor{baseline}{HTML}{84A8A1}

\begin{document}
% AI skills for reliable software
\begingroup
\setbeamercolor{background canvas}{bg=dark}
\begin{frame}[plain]
\note{\tiny
The comparison uses archived AI-generated outputs from November 18, 2026 and a separate comparison of C++ inference tests. The paired runs hold catalog access constant. The examples demonstrate concrete review
and testing benefits, with clear limits on causality.

}
\begin{slidecanvas}
\SlideHeader{AI skills for reliable software}{ENGINEERING PRESENTATION}{mint}{light}
\SlideFooter{MATHEMATICAL ABSTRACTIONS}{mint}{darkrule}
\SlideText{0.75}{2.02}{10.9}{14.5}{light}{}{From mathematical rules to code you can\par review and tests that expose defects}
\SlidePanel{.75}{4.39}{.34}{.04}{mint}
\SlideText{.75}{4.62}{9.5}{11}{mint}{\bfseries}{What changes when an AI agent gets repository skills as well as the catalog?}
\end{slidecanvas}
\end{frame}
\endgroup

% Roadmap
\begingroup
\setbeamercolor{background canvas}{bg=paper}
\begin{frame}[plain]
\note{\tiny
Introduce specification-derived testing. Compare the generated scheduler and wet-grass
tests, then recap the evidence and apply the workflow to one API.
}
\begin{slidecanvas}
\SlideHeader{From a catalog rule to evidence}{ROADMAP}{teal}{ink}
\SlideText{.74}{1.79}{11.8}{10.5}{muted}{}{Start with the method, then see how it changes test design and exposes implementation faults.}
\SlidePanel{.74}{2.55}{11.86}{.59}{panel}
\SlideText{.94}{2.68}{3.34}{9.5}{teal}{\bfseries}{1 · Foundations}
\SlideText{4.35}{2.68}{8}{9.5}{ink}{}{See how a specification gives a test its expected result.}
\SlidePanel{.74}{3.25}{11.86}{.59}{panel}
\SlideText{.94}{3.38}{3.34}{9.5}{teal}{\bfseries}{2 · Scheduler}
\SlideText{4.35}{3.38}{8}{9.5}{ink}{}{Compare generated tests for an expiring job scheduler.}
\SlidePanel{.74}{3.95}{11.86}{.59}{panel}
\SlideText{.94}{4.08}{3.34}{9.5}{teal}{\bfseries}{3 · Wet grass}
\SlideText{4.35}{4.08}{8}{9.5}{ink}{}{See which tests catch a wrong probability.}
\SlidePanel{.74}{4.65}{11.86}{.59}{panel}
\SlideText{.94}{4.78}{3.34}{9.5}{teal}{\bfseries}{4 · Takeaways}
\SlideText{4.35}{4.78}{8}{9.5}{ink}{}{Weigh the results and apply the method to one API.}
\SlideFooter{MATHEMATICAL ABSTRACTIONS}{muted}{rule}
\end{slidecanvas}
\end{frame}
\endgroup

\section{Foundations}
% Why derive tests from a specification?
\begingroup
\setbeamercolor{background canvas}{bg=paper}
\begin{frame}[plain]
\note{\tiny
Here formal testing means deriving executable checks from precisely stated behavior. A test can
falsify a claim with a concrete counterexample. Passing tests support the cases and input ranges
exercised. A proof has separate obligations and evidence. The scheduler case that follows
shows how two generated suites turn the same public contract into runnable tests.

}
\begin{slidecanvas}
\SlideHeader{Why derive tests from a specification?}{FORMAL TESTING}{teal}{ink}
\SlideText{0.74}{1.79}{11.8}{10.5}{muted}{}{An explicit rule tells us what should happen even when the implementation changes.}
\SlideFooter{Source  •  Repository mathematical-tests skill · operation and claim mapping}{muted}{rule}
\SlideCard{0.73}{2.83}{3.8}{2.65}{Expose hidden cases}{State which inputs are valid, what equality means and which outputs are allowed.}
\SlideCard{4.77}{2.83}{3.8}{2.65}{Predict the answer}{Compute it from a rule, a small model or a known result. Use that prediction to test the code.}
\SlideCard{8.81}{2.83}{3.8}{2.65}{Catch regressions}{Choose inputs that distinguish correct behavior from a plausible mistake. Keep the failing example.}
\SlideText{0.75}{5.45}{11.8}{9.5}{teal}{\bfseries}{A failing assertion gives a counterexample. Passing checks give evidence for the cases run.}
\end{slidecanvas}
\end{frame}
\endgroup

\section{Scheduler}
% Scheduler task and public API
\begingroup
\setbeamercolor{background canvas}{bg=paper}
\begin{frame}[plain]
\begin{slidecanvas}
\SlideHeader{Test a scheduler with expiring jobs}{ORIGINAL TASK · PYTHON UNIT TESTS}{teal}{ink}
\SlideText{.74}{1.79}{11.8}{10.5}{muted}{}{Write tests for a supplied scheduler through its public API. Each job has an ID, an integer priority, and an expiry time.}
\SlideCard{.74}{2.72}{3.8}{2.36}{\texttt{put}}{Add a job or replace one with the same ID. The contract places no capacity limit on live jobs.}
\SlideCard{4.77}{2.72}{3.8}{2.36}{\texttt{snapshot}}{Return all live jobs sorted by ID. Leave those jobs in the scheduler.}
\SlideCard{8.80}{2.72}{3.8}{2.36}{\texttt{take}}{Remove a live job with the lowest priority number. Any tied minimum is allowed.}
\SlideText{.75}{5.53}{11.8}{9.5}{teal}{\bfseries}{A job expiring at time 5 appears in a snapshot at time 4 and is gone at time 5.}
\SlideText{.75}{6.13}{11.8}{8.3}{ink}{}{Each timed call purges expired jobs. Time may stay the same or advance. A backwards call raises an error and preserves state.}
\SlideFooter{Source  •  Archived case 11 · TASK.md  ·  Full prompt in the comparison slide notes}{muted}{rule}
\end{slidecanvas}
\end{frame}
\endgroup

% The guided test catches silent job loss
\begingroup
\setbeamercolor{background canvas}{bg=paper}
\begin{frame}[plain]
\note{\tiny
The guided suite inserted 513 distinct live jobs at time zero, checked the complete snapshot and
drained them by priority. The baseline suite used independent histories and checked legal tie
choices, expiration and clone behavior. The extra injected fault ignored live insertions after
eight active IDs. This was selected after inspecting the suites, so the primary score remains
10/10 for both.
\par\medskip\textbf{Original task prompt}\par

\# Design a test suite for an opaque expiring job scheduler

Write standard-library `unittest` tests in `test\_scheduler.py` for `import scheduler\_api as
api`. A reference implementation is supplied as `scheduler\_api.py`; do not modify it. Other
correct implementations may use different private state types and choose different tied jobs.

Contract:

- `new()` creates an empty mutable scheduler with logical time 0. - `put(state, job\_id, priority,
expires\_at, now)` first advances time and purges every job with `expires\_at <= now`. It then
replaces any job with this ID. If the supplied expiration is already `<= now`, the result has no
job with that ID. The return value is `None`. - `take(state, now)` advances time, purges expired
jobs, removes and returns a live minimum-priority job as `(job\_id, priority, expires\_at)`. Any
tied minimum may be chosen. Empty `take` raises `KeyError`; advancing time and expiration still
take effect. - `snapshot(state, now)` advances time, purges expired jobs and returns a tuple of
live job rows sorted by job ID. It does not remove live jobs. - `clone(state)` returns an
independent scheduler with the same jobs and logical time. Changes to either copy must not affect
the other. - Time starts at 0 and may stay equal or increase. Every timed operation with `now`
less than that scheduler's last accepted time raises `ValueError` and leaves jobs and time
unchanged. No expiration or partial update may occur on this rejected call. A clone inherits its
source's last accepted time.

Job IDs are nonempty strings, priorities are arbitrary signed Python integers, expiration and
`now` are arbitrary Python integers. Inputs are of these types. Negative expiration simply denotes
an already expired job; negative `now` is backwards from the initial clock. No capacity limit
applies. Only these module functions are public; do not inspect private attributes or depend on
internal containers.

The suite should give a maintainer confidence in replacement, expiration, time errors, cloning,
mixed operation histories, and allowed tie choices—not just happy paths. Put a concise explanation
of each test group's purpose and oracle in `test-plan.md`, with a reproducible command and actual
results in `answer.md`. Run the suite. Do not repair production code, alter the contract or change
the catalog. The amount of code or number of tests is not itself a quality goal.

}
\begin{slidecanvas}
\SlideHeader{Two generated tests check the snapshot}{GENERATED PYTHON · SIDE BY SIDE}{teal}{ink}
\SlideText{.74}{1.79}{11.8}{10.1}{muted}{}{Both tests insert live jobs at time zero and compare the snapshot with the sorted input.}
\SlideCode{.73}{2.48}{5.82}{2.89}{WITHOUT SKILLS · SIX JOBS}{6.35}{state~=~api.new()\par \#~...~rows~contains~six~live~jobs~...\par for~row~in~rows:\par ~~~~self.assertIsNone(api.put(state,~*row,~0))\par expected~=~tuple(sorted(rows))\par for~\_~in~range(3):\par ~~~~result~=~api.snapshot(state,~0)\par ~~~~\#~...~tuple~shape~assertions~...\par ~~~~self.assertEqual(result,~expected)}
\SlideCode{6.79}{2.48}{5.82}{2.89}{WITH SKILLS · 513 JOBS}{6.35}{s~=~api.new()\par rows~=~[(\mbox{f'job-\{i:04d\}'},~(i~*~37)~\%~101~-~50,\par ~~~~~~~~~1000~+~i)~for~i~in~range(513)]\par for~row~in~rows:\par ~~~~self.put(s,~*row,~0)\par self.assertEqual(api.snapshot(s,~0),\par ~~~~~~~~~~~~~~~~~tuple(sorted(rows)))}
\SlideText{.75}{5.57}{11.8}{9.3}{teal}{\bfseries}{An extra fault capped the queue at eight jobs. Only the guided suite caught it.}
\SlideText{.75}{6.20}{11.8}{8.1}{ink}{}{Both complete suites caught all 10 predeclared faults. The eight-job cap was tried after reviewing the suites.}
\SlideFooter{Source  •  Archived case 11 · Abridged test\_scheduler.py excerpts and capacity-supplement.json}{muted}{rule}
\end{slidecanvas}
\end{frame}
\endgroup

\section{Wet grass}
\begingroup
\setbeamercolor{background canvas}{bg=paper}
\begin{frame}[plain]
\note{\tiny
The task was to write GoogleTest tests for the fixed C++ Bayesian network. The documented
contract covers joint probabilities, evidence matching, marginal probabilities and posterior
queries. Each of the four variables is Boolean. Evidence records true, false or unobserved
for each variable. The API throws std::domain\_error for zero-mass evidence, including
when the queried variable is already observed. The probability tables are illustrative.
}
\begin{slidecanvas}
\SlideHeader{The grass is wet. Did it rain?}{ORIGINAL TASK · TEST A C++ INFERENCE ENGINE}{teal}{ink}
\SlideText{.74}{1.79}{11.8}{10.5}{muted}{}{Given observations about cloud cover, the sprinkler or wet grass, the API computes the probability of rain.}
\SlideCard{.74}{2.64}{5.76}{2.52}{Wet grass alone}{Rain is about \textbf{71\%} likely in this model. Both suites check the exact reference ratio $0.4581/0.6471$.}
\SlideCard{6.79}{2.64}{5.76}{2.52}{Wet grass and sprinkler on}{Rain drops to about \textbf{32\%}. The running sprinkler gives another explanation for the wet grass. Both suites test this too.}
\SlideText{.75}{5.64}{11.8}{9.5}{teal}{\bfseries}{What happens when the grass is wet, the sprinkler is off and the sky is clear?}
\SlideFooter{Source  •  automatic\_tests/bayesian\_network.md and Bayesian network GoogleTest suites}{muted}{rule}
\end{slidecanvas}
\end{frame}

\begin{frame}[plain]
\note{\tiny
The baseline checks the complete joint distribution using independent expected numbers
and tests selected posterior cases. Both snippets are abridged from the original suites.
The guided oracle stores 16 documented joint probabilities as integer weights over 20,000
and selects compatible assignments using bit masks. No production helper computes its
expected values. The network factorization gives the weight table. The conditional-probability
ratio gives an expected result for each evidence/query pair. The implementation extends the
catalog network interface to already-observed query variables using event intersection.
Of 81 evidence patterns, 78 have positive mass. With 4 variables and 2 values this gives
624 posterior assertions against the oracle. All 8 queries on each of the 3 zero-mass
patterns must throw under the C++ contract.
}
\begin{slidecanvas}
\SlideHeader{From two weather examples to 624 queries}{GENERATED GOOGLETEST · WITHOUT AND WITH SKILLS}{teal}{ink}
\SlideText{.74}{1.78}{11.8}{10.1}{muted}{}{The catalog gives the test oracle its rule. For positive-mass evidence $B$, $P(A\mid B)=P(A\cap B)/P(B)$.}
\SlideCode{.73}{2.63}{5.82}{2.48}{WITHOUT SKILLS · SELECTED POSTERIORS}{6.7}{EXPECT\_NEAR(rain\_given\_wet,\par ~~~~0.4581~/~0.6471,~kTolerance);\par\medskip EXPECT\_NEAR(rain\_given\_wet\_and\_sprinkler,\par ~~~~0.0891~/~0.2781,~kTolerance);}
\SlideCode{6.79}{2.63}{5.82}{2.48}{WITH SKILLS · EXHAUSTIVE ORACLE}{6.7}{const~int~denominator~=~event\_weight(evidence);\par //~...~all~four~variables~and~both~values~...\par const~double~result~=~posterior(\par ~~~~variables[q],~value,~evidence);\par EXPECT\_NEAR(result,~double(numerator)~/\par ~~~~denominator,~tolerance);}
\SlideText{.75}{5.56}{11.8}{9}{teal}{\bfseries}{The guided suite enumerates all $3^4=81$ evidence patterns using true, false and unobserved.}
\SlideText{.75}{6.13}{11.8}{8.1}{ink}{}{Integer weights give an independent expected answer for each positive-mass query. Three impossible patterns exercise the API's exception path.}
\SlideFooter{Sources  •  ConditionalProbability-ratio  ·  automatic\_tests/\{noskill\_tests,tests\}/bayesian\_network\_test.cpp}{muted}{rule}
\end{slidecanvas}
\end{frame}

\begin{frame}[plain]
\note{\tiny
Local rerun on 2026-11-25. Compile the two suites separately against the unchanged source
using C++20 and GoogleTest. Then repeat with the same scratch change in each copy.
In posterior, immediately before its final return, count engaged evidence optionals
using has\_value(). If the count is variable\_count minus one, return
evidence\_probability(combined) directly. Leave the existing final return that divides
by denominator in place for other counts. The guided failures were
EveryPositiveEvidenceQueryMatchesExactRatio and ConditionalRowsAreNormalizedAndBounded.
For Cloudy=false, Sprinkler=false, WetGrass=true, the only positive-mass assignment
has Rain=true and weight 900/20,000. The faulty return is 0.045 and the correct
ratio is 0.045/0.045. Both suites pass 14/14 on the unchanged code. This is one
seeded fault in an existing pair of suites. The matched runs elsewhere in the deck
used a separate evaluation protocol.
}
\begin{slidecanvas}
\SlideHeader{One missed division changes the answer}{DELIBERATE FAULT · SAME C++ IMPLEMENTATION}{teal}{ink}
\SlideText{.74}{1.78}{11.8}{10.5}{muted}{}{Wet grass, sprinkler off, clear sky. In this model, rain is the only possible explanation.}
\SlideCard{.74}{2.55}{5.76}{1.54}{Correct answer · 100\%}{$P(\text{rain and evidence})/P(\text{evidence})=0.045/0.045$.}
\SlideCard{6.79}{2.55}{5.76}{1.54}{Faulty answer · 4.5\%}{With three observations, the seeded bug returns the joint mass directly.}
\SlideText{.76}{4.58}{4.2}{9}{teal}{\bfseries}{Test suite}
\SlideText{5.34}{4.58}{3.3}{9}{teal}{\bfseries}{Original code}
\SlideText{8.9}{4.58}{3.3}{9}{teal}{\bfseries}{With seeded fault}
\SlidePanel{.74}{4.98}{11.8}{.017}{rule}
\SlideText{.76}{5.17}{4.2}{9}{ink}{}{Without skills}
\SlideText{5.34}{5.17}{3.3}{9}{ink}{}{14/14 passed}
\SlideText{8.9}{5.17}{3.3}{9}{ink}{}{14/14 passed}
\SlideText{.76}{5.74}{4.2}{9}{ink}{}{With skills}
\SlideText{5.34}{5.74}{3.3}{9}{ink}{}{14/14 passed}
\SlideText{8.9}{5.74}{3.3}{9}{teal}{\bfseries}{2 tests failed}
\SlideFooter{Local rerun  •  25 November 2026  •  One seeded fault on this fixed four-variable network}{muted}{rule}
\end{slidecanvas}
\end{frame}
\endgroup

\section{Takeaways}
% What changed with skills
\begingroup
\setbeamercolor{background canvas}{bg=paper}
\begin{frame}[plain]
\note{\tiny
The scheduler comes from the matched pilot. Each arm had the same task and catalog,
with skill instructions exposed only in the guided arm. Both scheduler suites accepted both
correct implementations. The eight-job cap was picked after inspecting the suites.
Wet grass compares existing GoogleTest suites on an unchanged control and one deliberate
fault, outside the matched pilot. Both suites passed 14/14 on the control. These are
individual comparisons, with no estimate of average gains.
}
\begin{slidecanvas}
\SlideHeader{What changed with skills?}{EVIDENCE AND LIMITS}{teal}{ink}
\SlideText{.74}{1.79}{11.8}{10.5}{muted}{}{The scheduler suites tied on the predeclared faults. Two additional probes show where the guided tests went further.}
\SlideText{.98}{2.61}{6.2}{7}{teal}{\bfseries}{COMPARISON}
\SlideText{7.55}{2.61}{2.12}{7}{teal}{\bfseries}{WITHOUT SKILLS}
\SlideText{10.1}{2.61}{2.23}{7}{teal}{\bfseries}{WITH SKILLS}
\SlidePanel{.74}{2.98}{11.87}{.58}{panel}
\SlideText{.98}{3.09}{6.46}{8.7}{ink}{}{Matched · Correct schedulers accepted}
\SlideText{7.55}{3.09}{2.12}{8.7}{ink}{\bfseries}{2/2}
\SlideText{10.1}{3.09}{2.23}{8.7}{ink}{\bfseries}{2/2}
\SlidePanel{.74}{3.62}{11.87}{.58}{alternate}
\SlideText{.98}{3.73}{6.46}{8.7}{ink}{}{Matched · Predeclared scheduler faults caught}
\SlideText{7.55}{3.73}{2.12}{8.7}{ink}{\bfseries}{10/10}
\SlideText{10.1}{3.73}{2.23}{8.7}{ink}{\bfseries}{10/10}
\SlidePanel{.74}{4.26}{11.87}{.58}{panel}
\SlideText{.98}{4.37}{6.46}{8.7}{ink}{}{Post-hoc scheduler probe · Eight-job cap}
\SlideText{7.55}{4.37}{2.12}{8.7}{ink}{}{Missed}
\SlideText{10.1}{4.37}{2.23}{8.7}{teal}{\bfseries}{Caught}
\SlidePanel{.74}{4.90}{11.87}{.58}{alternate}
\SlideText{.98}{5.01}{6.46}{8.7}{ink}{}{Separate wet-grass probe · Missing normalization}
\SlideText{7.55}{5.01}{2.12}{8.7}{ink}{}{Missed}
\SlideText{10.1}{5.01}{2.23}{8.7}{teal}{\bfseries}{Caught}
\SlideText{.75}{5.91}{11.8}{9}{ink}{}{The scheduler test varied queue size. The wet-grass test varied evidence combinations. Both probes show how broader inputs can expose faults.}
\SlideFooter{Sources  •  Archived matched pilot and capacity supplement  ·  Separate Bayesian network rerun}{muted}{rule}
\end{slidecanvas}
\end{frame}
\endgroup

% Put the skills to work on one risky API
\begingroup
\setbeamercolor{background canvas}{bg=paper}
\begin{frame}[plain]
\note{\tiny
Example prompt for the presenter: Read ai/skills/mathematical-tests/SKILL.md.
Map the rule that jobs expire at or before the current time to the public
scheduler API. Insert one job that expires at time 5 and check that it appears
at time 4 and is gone at time 5. Run the generated test against correct code
and a temporary copy that expires jobs one tick late. Keep the failing assertion.
}
\begin{slidecanvas}
\SlideHeader{Put the skills to work on one risky API}{ADOPTION}{teal}{ink}
\SlideCard{.73}{2.58}{3.8}{2.25}{Pick a rule}{Choose a behavior your API must guarantee.}
\SlideCard{4.77}{2.58}{3.8}{2.25}{Generate a test}{Use the skill to turn the rule into an assertion.}
\SlideCard{8.81}{2.58}{3.8}{2.25}{Check for defects}{Run the test against correct code and a deliberate bug.}
\SlideFooter{Repository skill  •  Mathematical tests}{muted}{rule}
\end{slidecanvas}
\end{frame}
\endgroup

% Source artifacts and reproduction
\begingroup
\setbeamercolor{background canvas}{bg=paper}
\begin{frame}[plain]
\note{\tiny
Slides 4--5 use Git commit a71a69d, under
ai/evaluations/results/2026-11-18-hard-quality/. Use git archive a71a69d with that
path to extract it into a scratch directory. Case 11 contains test\_scheduler.py
and test-plan.md in the with-skills and without-skills directories.
For each arm, run \texttt{python3 ai/evaluations/quality\_checks.py -{}-case 11 -{}-workspace PATH}.
PATH is the extracted arm directory. Read capacity-supplement.json for the
separate eight-job fault, source and logs.

\par\medskip
Slides 6--8 use the sibling automatic\_tests workspace. Compile each
bayesian\_network\_test.cpp with C++20 and GoogleTest, preserving the relative
include of bayesian\_network.cpp. The suites are under noskill\_tests/ and tests/.
Run each against the original target. Repeat against a scratch copy with the edit
specified in slide 8's notes. The integer-weight oracle and catalog mapping are
in tests/bayesian\_network\_tests.md.
}
\begin{slidecanvas}
\SlideHeading{Where to verify each example}{APPENDIX}{teal}{ink}
\SlideText{.74}{1.79}{11.8}{9.5}{muted}{}{Scheduler archive · Git commit a71a69d · ai/evaluations/results/2026-11-18-hard-quality/}
\SlideText{.75}{2.72}{11.7}{9.5}{teal}{\bfseries}{Slides 4–5 · Scheduler task, tests and extra fault}
\SlideText{.75}{3.12}{11.7}{8.4}{ink}{}{Archive case-11/ · Both test\_scheduler.py files. Eight-job fault · capacity-supplement.json.}
\SlideText{.75}{4.22}{11.7}{9.5}{teal}{\bfseries}{Slides 6–8 · Wet-grass tests and seeded fault}
\SlideText{.75}{4.62}{11.7}{8.4}{ink}{}{Companion automatic\_tests workspace · Compare bayesian\_network\_test.cpp in noskill\_tests/ and tests/.\par Read tests/bayesian\_network\_tests.md for the oracle and slide 8's notes for the exact fault edit.}
\SlideFooter{Rerun instructions are in the presenter notes.}{muted}{rule}
\end{slidecanvas}
\end{frame}
\endgroup

\end{document}
